% Options for packages loaded elsewhere
% Options for packages loaded elsewhere
\PassOptionsToPackage{unicode}{hyperref}
\PassOptionsToPackage{hyphens}{url}
\PassOptionsToPackage{dvipsnames,svgnames,x11names}{xcolor}
%
\documentclass[
  spanish,
  10.5pt,
  a4paper,
  DIV=11,
  numbers=noendperiod]{scrartcl}
\usepackage{xcolor}
\usepackage[top=24mm,bottom=24mm,left=23mm,right=23mm]{geometry}
\usepackage{amsmath,amssymb}
\setcounter{secnumdepth}{5}
\usepackage{iftex}
\ifPDFTeX
  \usepackage[T1]{fontenc}
  \usepackage[utf8]{inputenc}
  \usepackage{textcomp} % provide euro and other symbols
\else % if luatex or xetex
  \usepackage{unicode-math} % this also loads fontspec
  \defaultfontfeatures{Scale=MatchLowercase}
  \defaultfontfeatures[\rmfamily]{Ligatures=TeX,Scale=1}
\fi
\usepackage{lmodern}
\ifPDFTeX\else
  % xetex/luatex font selection
  \setmainfont[]{Latin Modern Roman}
  \setsansfont[]{Latin Modern Sans}
  \setmonofont[]{Latin Modern Mono}
\fi
% Use upquote if available, for straight quotes in verbatim environments
\IfFileExists{upquote.sty}{\usepackage{upquote}}{}
\IfFileExists{microtype.sty}{% use microtype if available
  \usepackage[]{microtype}
  \UseMicrotypeSet[protrusion]{basicmath} % disable protrusion for tt fonts
}{}
\makeatletter
\@ifundefined{KOMAClassName}{% if non-KOMA class
  \IfFileExists{parskip.sty}{%
    \usepackage{parskip}
  }{% else
    \setlength{\parindent}{0pt}
    \setlength{\parskip}{6pt plus 2pt minus 1pt}}
}{% if KOMA class
  \KOMAoptions{parskip=half}}
\makeatother
% Make \paragraph and \subparagraph free-standing
\makeatletter
\ifx\paragraph\undefined\else
  \let\oldparagraph\paragraph
  \renewcommand{\paragraph}{
    \@ifstar
      \xxxParagraphStar
      \xxxParagraphNoStar
  }
  \newcommand{\xxxParagraphStar}[1]{\oldparagraph*{#1}\mbox{}}
  \newcommand{\xxxParagraphNoStar}[1]{\oldparagraph{#1}\mbox{}}
\fi
\ifx\subparagraph\undefined\else
  \let\oldsubparagraph\subparagraph
  \renewcommand{\subparagraph}{
    \@ifstar
      \xxxSubParagraphStar
      \xxxSubParagraphNoStar
  }
  \newcommand{\xxxSubParagraphStar}[1]{\oldsubparagraph*{#1}\mbox{}}
  \newcommand{\xxxSubParagraphNoStar}[1]{\oldsubparagraph{#1}\mbox{}}
\fi
\makeatother

\usepackage{color}
\usepackage{fancyvrb}
\newcommand{\VerbBar}{|}
\newcommand{\VERB}{\Verb[commandchars=\\\{\}]}
\DefineVerbatimEnvironment{Highlighting}{Verbatim}{commandchars=\\\{\}}
% Add ',fontsize=\small' for more characters per line
\usepackage{framed}
\definecolor{shadecolor}{RGB}{241,243,245}
\newenvironment{Shaded}{\begin{snugshade}}{\end{snugshade}}
\newcommand{\AlertTok}[1]{\textcolor[rgb]{0.68,0.00,0.00}{#1}}
\newcommand{\AnnotationTok}[1]{\textcolor[rgb]{0.37,0.37,0.37}{#1}}
\newcommand{\AttributeTok}[1]{\textcolor[rgb]{0.40,0.45,0.13}{#1}}
\newcommand{\BaseNTok}[1]{\textcolor[rgb]{0.68,0.00,0.00}{#1}}
\newcommand{\BuiltInTok}[1]{\textcolor[rgb]{0.00,0.23,0.31}{#1}}
\newcommand{\CharTok}[1]{\textcolor[rgb]{0.13,0.47,0.30}{#1}}
\newcommand{\CommentTok}[1]{\textcolor[rgb]{0.37,0.37,0.37}{#1}}
\newcommand{\CommentVarTok}[1]{\textcolor[rgb]{0.37,0.37,0.37}{\textit{#1}}}
\newcommand{\ConstantTok}[1]{\textcolor[rgb]{0.56,0.35,0.01}{#1}}
\newcommand{\ControlFlowTok}[1]{\textcolor[rgb]{0.00,0.23,0.31}{\textbf{#1}}}
\newcommand{\DataTypeTok}[1]{\textcolor[rgb]{0.68,0.00,0.00}{#1}}
\newcommand{\DecValTok}[1]{\textcolor[rgb]{0.68,0.00,0.00}{#1}}
\newcommand{\DocumentationTok}[1]{\textcolor[rgb]{0.37,0.37,0.37}{\textit{#1}}}
\newcommand{\ErrorTok}[1]{\textcolor[rgb]{0.68,0.00,0.00}{#1}}
\newcommand{\ExtensionTok}[1]{\textcolor[rgb]{0.00,0.23,0.31}{#1}}
\newcommand{\FloatTok}[1]{\textcolor[rgb]{0.68,0.00,0.00}{#1}}
\newcommand{\FunctionTok}[1]{\textcolor[rgb]{0.28,0.35,0.67}{#1}}
\newcommand{\ImportTok}[1]{\textcolor[rgb]{0.00,0.46,0.62}{#1}}
\newcommand{\InformationTok}[1]{\textcolor[rgb]{0.37,0.37,0.37}{#1}}
\newcommand{\KeywordTok}[1]{\textcolor[rgb]{0.00,0.23,0.31}{\textbf{#1}}}
\newcommand{\NormalTok}[1]{\textcolor[rgb]{0.00,0.23,0.31}{#1}}
\newcommand{\OperatorTok}[1]{\textcolor[rgb]{0.37,0.37,0.37}{#1}}
\newcommand{\OtherTok}[1]{\textcolor[rgb]{0.00,0.23,0.31}{#1}}
\newcommand{\PreprocessorTok}[1]{\textcolor[rgb]{0.68,0.00,0.00}{#1}}
\newcommand{\RegionMarkerTok}[1]{\textcolor[rgb]{0.00,0.23,0.31}{#1}}
\newcommand{\SpecialCharTok}[1]{\textcolor[rgb]{0.37,0.37,0.37}{#1}}
\newcommand{\SpecialStringTok}[1]{\textcolor[rgb]{0.13,0.47,0.30}{#1}}
\newcommand{\StringTok}[1]{\textcolor[rgb]{0.13,0.47,0.30}{#1}}
\newcommand{\VariableTok}[1]{\textcolor[rgb]{0.07,0.07,0.07}{#1}}
\newcommand{\VerbatimStringTok}[1]{\textcolor[rgb]{0.13,0.47,0.30}{#1}}
\newcommand{\WarningTok}[1]{\textcolor[rgb]{0.37,0.37,0.37}{\textit{#1}}}

\usepackage{longtable,booktabs,array}
\newcounter{none} % for unnumbered tables
\usepackage{calc} % for calculating minipage widths
% Correct order of tables after \paragraph or \subparagraph
\usepackage{etoolbox}
\makeatletter
\patchcmd\longtable{\par}{\if@noskipsec\mbox{}\fi\par}{}{}
\makeatother
% Allow footnotes in longtable head/foot
\IfFileExists{footnotehyper.sty}{\usepackage{footnotehyper}}{\usepackage{footnote}}
\makesavenoteenv{longtable}
\usepackage{graphicx}
\makeatletter
\newsavebox\pandoc@box
\newcommand*\pandocbounded[1]{% scales image to fit in text height/width
  \sbox\pandoc@box{#1}%
  \Gscale@div\@tempa{\textheight}{\dimexpr\ht\pandoc@box+\dp\pandoc@box\relax}%
  \Gscale@div\@tempb{\linewidth}{\wd\pandoc@box}%
  \ifdim\@tempb\p@<\@tempa\p@\let\@tempa\@tempb\fi% select the smaller of both
  \ifdim\@tempa\p@<\p@\scalebox{\@tempa}{\usebox\pandoc@box}%
  \else\usebox{\pandoc@box}%
  \fi%
}
% Set default figure placement to htbp
\def\fps@figure{htbp}
\makeatother



\ifLuaTeX
\usepackage[bidi=basic,shorthands=off]{babel}
\else
\usepackage[bidi=default,shorthands=off]{babel}
\fi
\ifPDFTeX
\else
\babelfont{rm}[]{Latin Modern Roman}
\fi
\ifLuaTeX
  \usepackage{selnolig} % disable illegal ligatures
\fi


\setlength{\emergencystretch}{3em} % prevent overfull lines

\providecommand{\tightlist}{%
  \setlength{\itemsep}{0pt}\setlength{\parskip}{0pt}}



 


\usepackage{fvextra}
\DefineVerbatimEnvironment{Highlighting}{Verbatim}{commandchars=\\\{\},breaklines=true,fontsize=\footnotesize}
\usepackage{scrlayer-scrpage}
\clearpairofpagestyles
\ihead{\small SBVI · Ingeniería Biomédica}
\ohead{\small EMGDB · Datos reales}
\cfoot{\pagemark}
\setkomafont{pageheadfoot}{\normalfont\sffamily\footnotesize}
\usepackage{microtype}
\usepackage{needspace}
\usepackage{tikz}
\usepackage{etoolbox}
\pretocmd{\section}{\Needspace{14\baselineskip}}{}{}
\setlength{\emergencystretch}{3em}
\setlength{\parskip}{5pt}
\setcounter{tocdepth}{1}
\KOMAoption{captions}{tableheading}
\makeatletter
\@ifpackageloaded{caption}{}{\usepackage{caption}}
\AtBeginDocument{%
\ifdefined\contentsname
  \renewcommand*\contentsname{Tabla de contenidos}
\else
  \newcommand\contentsname{Tabla de contenidos}
\fi
\ifdefined\listfigurename
  \renewcommand*\listfigurename{Listado de Figuras}
\else
  \newcommand\listfigurename{Listado de Figuras}
\fi
\ifdefined\listtablename
  \renewcommand*\listtablename{Listado de Tablas}
\else
  \newcommand\listtablename{Listado de Tablas}
\fi
\ifdefined\figurename
  \renewcommand*\figurename{Figura}
\else
  \newcommand\figurename{Figura}
\fi
\ifdefined\tablename
  \renewcommand*\tablename{Tabla}
\else
  \newcommand\tablename{Tabla}
\fi
}
\@ifpackageloaded{float}{}{\usepackage{float}}
\floatstyle{ruled}
\@ifundefined{c@chapter}{\newfloat{codelisting}{h}{lop}}{\newfloat{codelisting}{h}{lop}[chapter]}
\floatname{codelisting}{Listado}
\newcommand*\listoflistings{\listof{codelisting}{Listado de Listados}}
\makeatother
\makeatletter
\makeatother
\makeatletter
\@ifpackageloaded{caption}{}{\usepackage{caption}}
\@ifpackageloaded{subcaption}{}{\usepackage{subcaption}}
\makeatother
\usepackage{bookmark}
\IfFileExists{xurl.sty}{\usepackage{xurl}}{} % add URL line breaks if available
\urlstyle{same}
\hypersetup{
  pdftitle={Dashboard de EMG con datos abiertos},
  pdfauthor={Señales Bioeléctricas y Visualización Interactiva},
  pdflang={es},
  colorlinks=true,
  linkcolor={teal},
  filecolor={Maroon},
  citecolor={Blue},
  urlcolor={teal},
  pdfcreator={LaTeX via pandoc}}


\title{Dashboard de EMG con datos abiertos}
\usepackage{etoolbox}
\makeatletter
\providecommand{\subtitle}[1]{% add subtitle to \maketitle
  \apptocmd{\@title}{\par {\large #1 \par}}{}{}
}
\makeatother
\subtitle{Radiculopatía L5 · De la fisiología a la ingeniería}
\author{Señales Bioeléctricas y Visualización Interactiva}
\date{}
\begin{document}
\maketitle

\renewcommand*\contentsname{Tabla de contenidos}
{
\setcounter{tocdepth}{1}
\tableofcontents
}

\newpage

\section{Solicitud del médico y resultado
esperado}\label{solicitud-del-muxe9dico-y-resultado-esperado}

\begin{quote}
\textbf{Solicitud clínica ficticia para el equipo de ingeniería}\\
Atiendo personas con neuropatía asociada a radiculopatía L5. Necesito
una herramienta para revisar registros EMG del tibial anterior:
localizar segmentos, examinar su calidad y describir cómo cambian la
amplitud y el contenido espectral durante la exploración. Quiero
conservar la señal original junto a los indicadores y documentar
exactamente qué se calculó. En una etapa posterior me interesaría
comparar visitas, siempre que los registros fueran comparables. Para
desarrollar el prototipo, utilicen datos clínicos abiertos y no inventen
pacientes, mediciones ni cambios de tratamiento.
\end{quote}

\textbf{Producto.} Construirá un dashboard local en Python y Streamlit
sobre el dataset libre \textbf{Examples of Electromyograms, EMGDB 1.0.0,
de PhysioNet}. El caso principal será \texttt{emg\_neuropathy}; los
otros dos registros se utilizarán como ejemplos de contraste. La
aplicación descarga o lee los archivos originales y verifica su
integridad. No contiene un modo de pacientes simulados.

El PDF es una guía estática. Los controles interactivos funcionan en
\texttt{app.py}. El archivo Quarto incluye el código completo de los
módulos, las explicaciones, las fórmulas, las instrucciones de ejecución
y la evaluación.

\subsection{Alcance real del
monitoreo}\label{alcance-real-del-monitoreo}

Se implementa \textbf{exploración temporal dentro de un registro}. El
dataset contiene tres ejemplos de personas diferentes, no una serie
longitudinal. Por tanto, permite aprender a construir y verificar el
sistema, pero no demostrar mejoría, empeoramiento, respuesta terapéutica
ni capacidad diagnóstica poblacional.

Una etiqueta clínica suministrada por el repositorio es información de
contexto; no es una predicción producida por nuestro algoritmo. Los
descriptores de señal tampoco son equivalentes a la gravedad de la
lesión. No se usa un semáforo de ``normal/patológico'' ni un puntaje
compuesto de enfermedad.

\subsection{Articulación con la
asignatura}\label{articulaciuxf3n-con-la-asignatura}

La actividad desarrolla la secuencia del microcurrículo adjunto:
fisiología muscular, adquisición, análisis temporal y espectral,
visualización e integración. Aporta a RA1--RA5 y al módulo de EMG. Aquí
se analiza \textbf{EMG intramuscular}, por lo que no se traslada sin
revisión el procesamiento de EMG superficial.

\textbf{Prerrequisitos:} arreglos NumPy, archivos, frecuencia de
muestreo, filtros y Fourier. \textbf{Dedicación orientativa:} 8 a 10
horas. Se destinan 2 h al problema y los datos, 3 h al procesamiento, 2
h a la interfaz y 1 a 3 h a la verificación y exposición.

\textbf{Evidencias:} justificar cada indicador desde la fisiología,
comprobar las unidades desde la cabecera, reproducir cálculos sobre
datos reales y explicar los límites clínicos del tablero.

\section{Dataset libre y trazabilidad de la
fuente}\label{dataset-libre-y-trazabilidad-de-la-fuente}

\subsection{Acceso, licencia y
atribución}\label{acceso-licencia-y-atribuciuxf3n}

El dataset está disponible sin cuenta ni aprobación individual en
\href{https://physionet.org/content/emgdb/1.0.0/}{PhysioNet EMGDB
1.0.0}, DOI \href{https://doi.org/10.13026/C24S3D}{10.13026/C24S3D}. Su
licencia es \textbf{Open Data Commons Attribution License v1.0 (ODC-By
1.0)}. Mantenga la atribución, el vínculo a la licencia y la
identificación de las transformaciones al redistribuir un resultado
derivado {[}1{]}.

Los ejemplos fueron facilitados por Seward Rutkove, Beth Israel
Deaconess Medical Center/Harvard Medical School. Son registros de un
canal del tibial anterior. Las cabeceras documentan 4000 muestras/s y
una calibración de 10000 cuentas/mV {[}2--4{]}.

\Needspace{10\baselineskip}

{\def\LTcaptype{none} % do not increment counter
\begin{longtable}[]{@{}
  >{\raggedright\arraybackslash}p{(\linewidth - 6\tabcolsep) * \real{0.2500}}
  >{\raggedright\arraybackslash}p{(\linewidth - 6\tabcolsep) * \real{0.2500}}
  >{\raggedleft\arraybackslash}p{(\linewidth - 6\tabcolsep) * \real{0.2500}}
  >{\raggedleft\arraybackslash}p{(\linewidth - 6\tabcolsep) * \real{0.2500}}@{}}
\toprule\noalign{}
\begin{minipage}[b]{\linewidth}\raggedright
Registro original
\end{minipage} & \begin{minipage}[b]{\linewidth}\raggedright
Etiqueta publicada
\end{minipage} & \begin{minipage}[b]{\linewidth}\raggedleft
Muestras
\end{minipage} & \begin{minipage}[b]{\linewidth}\raggedleft
Duración calculada
\end{minipage} \\
\midrule\noalign{}
\endhead
\bottomrule\noalign{}
\endlastfoot
\texttt{emg\_neuropathy} & Neuropatía por radiculopatía L5 derecha &
147858 & 36.9645 s \\
\texttt{emg\_healthy} & Sin historia de enfermedad neuromuscular & 50860
& 12.7150 s \\
\texttt{emg\_myopathy} & Miopatía por polimiositis & 110337 & 27.58425
s \\
\end{longtable}
}

Las duraciones son \(N/f_s\), calculadas a partir de las cabeceras y
confirmadas al leer los binarios. No son tres fases de una enfermedad ni
tres condiciones de una misma persona.

\subsection{Qué información está disponible y cuál
falta}\label{quuxe9-informaciuxf3n-estuxe1-disponible-y-cuuxe1l-falta}

Los archivos contienen EMG registrada durante una tarea muscular y su
transición a relajación, según la descripción de origen. No incorporan
un canal de fuerza ni anotaciones temporales expertas que delimiten cada
potencial motor. El usuario no debe asignar automáticamente a toda
ventana una condición fisiológica solo por su amplitud.

La ficha indica adquisición inicial a 50 kHz y reducción posterior a 4
kHz. No se recuperan las frecuencias que se perdieron. Con 4 kHz,
Nyquist es 2 kHz, independientemente de la banda original del equipo. El
filtro digital posterior no corrige aliasing previo {[}1{]}.

El proyecto emplea archivos \texttt{.dat} y \texttt{.hea}. Se detectó
una inconsistencia de separación numérica en la copia \texttt{.txt} del
registro de neuropatía; no se corrige inventando números ni eliminando
filas. Los binarios descargados se verifican contra los SHA-256
publicados {[}5{]}.

\subsection{Archivos que utilizará}\label{archivos-que-utilizaruxe1}

Cada registro necesita un archivo binario y su cabecera. El paquete
incluye los seis archivos, el listado oficial de hashes y una nota de
atribución. Si trabaja solo con el QMD, el módulo \texttt{datos.py} los
descarga desde PhysioNet al ejecutarse.

El lector se limita deliberadamente a \textbf{estos archivos de un
canal, formato WFDB 16}. No es un lector general de todos los formatos
WFDB. Rechaza cambios de formato, ganancia, frecuencia, longitud o
integridad.

\section{De la lesión al indicador de
ingeniería}\label{de-la-lesiuxf3n-al-indicador-de-ingenieruxeda}

\subsection{Qué mide el electrodo}\label{quuxe9-mide-el-electrodo}

Una unidad motora agrupa una motoneurona y las fibras musculares que
inerva. La EMG registra contribuciones de potenciales de acción cercanos
al electrodo. Un modelo conceptual es

\[x(t)=\sum_{k=1}^{K}\sum_j a_k h_k(t-t_{kj})+\eta(t),\]

donde \(h_k\) es la forma del potencial observado de la unidad \(k\),
\(t_{kj}\) sus instantes de descarga, \(a_k\) recoge la geometría y
conducción, y \(\eta\) representa ruido y otras contribuciones. El
electrodo intramuscular observa un volumen diferente del de la EMG
superficial. Sus amplitudes y morfologías no se intercambian
directamente.

En una lesión de raíz que afecte axones motores puede alterarse el
conjunto de unidades activables. En procesos crónicos puede existir
reinervación colateral y remodelación de unidades. Estos fenómenos
pueden modificar la morfología y el patrón de interferencia, pero su
evaluación exige contexto, esfuerzo y muestreo muscular adecuados. Un
estudio de simulación de unidades reinervadas muestra precisamente que
las propiedades observadas dependen de la organización de fibras y de la
técnica de registro {[}6{]}.

\subsection{Separar tres niveles de
evidencia}\label{separar-tres-niveles-de-evidencia}

\textbf{Fenómeno fisiológico:} pérdida o remodelación de unidades
motoras. \textbf{Hallazgo electrofisiológico:} cambios en morfología o
reclutamiento evaluados por un especialista. \textbf{Descriptor
computacional:} un número calculado sobre una ventana de señal. Pasar
del tercer nivel al primero requiere una validación que este dataset no
proporciona.

El tutorial construye descriptores exploratorios útiles para organizar
la revisión. No estima número de unidades motoras, grado de denervación,
porcentaje de recuperación ni duración de un potencial aislado.

\Needspace{12\baselineskip}

{\def\LTcaptype{none} % do not increment counter
\begin{longtable}[]{@{}
  >{\raggedright\arraybackslash}p{(\linewidth - 4\tabcolsep) * \real{0.3333}}
  >{\raggedright\arraybackslash}p{(\linewidth - 4\tabcolsep) * \real{0.3333}}
  >{\raggedright\arraybackslash}p{(\linewidth - 4\tabcolsep) * \real{0.3333}}@{}}
\toprule\noalign{}
\begin{minipage}[b]{\linewidth}\raggedright
Pregunta fisiológica
\end{minipage} & \begin{minipage}[b]{\linewidth}\raggedright
Medida implementada
\end{minipage} & \begin{minipage}[b]{\linewidth}\raggedright
Límite de interpretación
\end{minipage} \\
\midrule\noalign{}
\endhead
\bottomrule\noalign{}
\endlastfoot
¿Cómo varía la amplitud eléctrica del segmento? & RMS de la EMG filtrada
y centrada & También cambia con esfuerzo y posición del electrodo \\
¿Qué excursión de amplitud contiene la ventana? & Máximo menos mínimo &
No equivale a amplitud de una unidad aislada \\
¿Predominan excursiones puntuales frente al nivel medio? & Factor de
cresta & Un artefacto también puede elevarlo \\
¿Dónde se distribuye la potencia en la banda analizada? & Frecuencia
mediana de 20--1000 Hz & No es frecuencia de descarga ni diagnóstico de
neuropatía \\
\end{longtable}
}

\subsection{Qué requeriría un indicador más
específico}\label{quuxe9-requeriruxeda-un-indicador-muxe1s-especuxedfico}

Para cuantificar duración, fases o amplitud de potenciales de unidad
motora se necesitarían eventos aislados o descompuestos, límites de
inicio y final validados y revisión experta. La anchura de cualquier
pico de la mezcla no cumple esa condición. Para estudiar reclutamiento
se necesita relacionar actividad con el nivel de activación o fuerza;
contar picos no cuenta unidades.

\textbf{Pregunta de verificación.} Si el RMS aumenta mientras cambia el
esfuerzo, ¿se puede atribuir el cambio a reinervación? Identifique el
factor de confusión y qué medición adicional haría falta antes de
formular esa conclusión.

\section{Diseño matemático y
temporal}\label{diseuxf1o-matemuxe1tico-y-temporal}

\subsection{Preprocesamiento
declarado}\label{preprocesamiento-declarado}

El original permanece disponible para inspección. La rama analítica
aplica un pasa banda Butterworth de 20--1000 Hz con orden de diseño 4 y
secciones de segundo orden. La banda es una \textbf{decisión docente},
no un estándar diagnóstico ni un sustituto del registro original.
Limitar el extremo superior elimina información que podría ser relevante
para morfología de potenciales; por eso no se usa la salida para medir
unidades aisladas.

El filtrado bidireccional \texttt{sosfiltfilt} utiliza muestras futuras:
es offline. La transformación pasa banda y la aplicación en dos sentidos
modifican el orden y la respuesta efectivos; no se describe como un
filtro causal de cuarto orden. Se filtra el registro continuo y después
se recortan ventanas. El ejemplo excluye 0.25 s por extremo; este margen
es una precaución de diseño que debe verificarse cuando cambien filtros
o datos.

No se rectifica antes de estimar el espectro. Aquí interesa el contenido
espectral de la EMG, no una frecuencia lenta de temblor. Tampoco se
aplica notch sin evidencia de interferencia y sin documentar su efecto.

\subsection{Amplitud y carácter
impulsivo}\label{amplitud-y-caruxe1cter-impulsivo}

Sea \(y[n]\) la señal filtrada dentro de una ventana de \(N\) muestras.
Defina \(z[n]=y[n]-\bar y\) para retirar su media. Las medidas son

\[\mathrm{RMS}=\sqrt{\frac{1}{N}\sum_{n=0}^{N-1}z[n]^2},\]

\[A_{pp}=\max_n z[n]-\min_n z[n],\qquad
C=\frac{\max_n|z[n]|}{\mathrm{RMS}}.\]

RMS y \(A_{pp}\) se expresan en microvoltios. \(C\) es adimensional y se
deja no disponible si el denominador es prácticamente nulo. El pico a
pico es sensible a extremos; al alargar la ventana aumenta la
posibilidad de incluir un pico grande. Por eso no deben compararse
ventanas de duración diferente sin reconocer ese efecto.

Si todas las amplitudes se multiplican por dos, RMS y \(A_{pp}\) se
duplican, mientras \(C\) permanece igual. Esta propiedad permite
verificar el código usando el registro real, sin fabricar pacientes ni
asumir un resultado clínico.

\subsection{Espectro y frecuencia
mediana}\label{espectro-y-frecuencia-mediana}

Se estima la densidad espectral \(S_z(f)\) con Welch sobre segmentos
Hann de 0.25 s y solapamiento del 50 \%. En 4000 Hz, cada segmento
contiene 1000 muestras, por lo que la separación entre bins es 4 Hz
{[}7{]}.

La frecuencia mediana restringida satisface

\[\int_{20}^{f_{med}}S_z(f)\,df
=\frac{1}{2}\int_{20}^{1000}S_z(f)\,df.\]

La implementación integra por trapecios e interpola el punto de mitad de
potencia. La PSD tiene unidades de \(\mu V^2/Hz\). La interpolación no
mejora la resolución física del estimador; imprimir muchos decimales no
aporta precisión biológica.

\(f_{med}\) puede cambiar por morfología, mezcla de unidades, filtros o
ruido. Un descenso no demuestra fatiga en este protocolo y una
diferencia entre registros no establece gravedad. No hay fuerza,
velocidad de conducción ni anotaciones suficientes para identificar el
mecanismo solo a partir de esa cifra.

\subsection{Escalas temporales}\label{escalas-temporales}

\Needspace{12\baselineskip}

{\def\LTcaptype{none} % do not increment counter
\begin{longtable}[]{@{}
  >{\raggedright\arraybackslash}p{(\linewidth - 4\tabcolsep) * \real{0.3333}}
  >{\raggedright\arraybackslash}p{(\linewidth - 4\tabcolsep) * \real{0.3333}}
  >{\raggedright\arraybackslash}p{(\linewidth - 4\tabcolsep) * \real{0.3333}}@{}}
\toprule\noalign{}
\begin{minipage}[b]{\linewidth}\raggedright
Escala
\end{minipage} & \begin{minipage}[b]{\linewidth}\raggedright
Configuración
\end{minipage} & \begin{minipage}[b]{\linewidth}\raggedright
Justificación
\end{minipage} \\
\midrule\noalign{}
\endhead
\bottomrule\noalign{}
\endlastfoot
Intervalo elegido & Inicialmente 0.25--5.25 s & Delimita la región que
se inspeccionará \\
Ventana analítica & 0.5, 1 o 2 s & Compromiso entre estabilidad y
localización \\
Paso & Mitad de la ventana & Tendencia más densa, con dependencia entre
puntos \\
Welch & 0.25 s, solapamiento 0.125 s & Grilla de 4 Hz; varios promedios
por ventana \\
Inspección de morfología & Zoom temporal del gráfico & Revisar eventos
sin asignarles identidad automática \\
\end{longtable}
}

Una ventana de 1 s tiene 4000 muestras y siete segmentos Welch
solapados; no son siete observaciones independientes. Los límites son
\([i,j)\): incluyen \(i\) y excluyen \(j\). El centro temporal etiqueta
el punto de la tendencia. Una transición brusca o un artefacto dentro de
la ventana limita su interpretación, aunque la fórmula produzca un
número.

\section{Preparar el proyecto y descargar los
datos}\label{preparar-el-proyecto-y-descargar-los-datos}

\subsection{Entorno reproducible}\label{entorno-reproducible}

Use Python 3.11 o 3.12, uv y Quarto CLI. En una carpeta nueva:

\begin{Shaded}
\begin{Highlighting}[]
\ExtensionTok{uv}\NormalTok{ init }\AttributeTok{{-}{-}python}\NormalTok{ 3.11}
\ExtensionTok{uv}\NormalTok{ add numpy scipy pandas streamlit plotly}
\ExtensionTok{uv}\NormalTok{ lock}
\end{Highlighting}
\end{Shaded}

Copie los bloques de las secciones de código en los archivos indicados.
Los fragmentos de un mismo archivo se concatenan en orden, separados por
una línea en blanco. El paquete completo ya contiene esos archivos y los
datos originales verificados.

\begin{Shaded}
\begin{Highlighting}[]
\ExtensionTok{uv}\NormalTok{ run python datos.py}
\ExtensionTok{uv}\NormalTok{ run python verificar.py}
\ExtensionTok{uv}\NormalTok{ run streamlit run app.py}
\end{Highlighting}
\end{Shaded}

\texttt{datos.py} obtiene los archivos que falten y comprueba sus
hashes. Una copia presente con hash incorrecto causa un error; no se
procesa ni se reemplaza silenciosamente. Con los archivos del paquete,
el análisis funciona sin conexión.

Para crear el PDF:

\begin{Shaded}
\begin{Highlighting}[]
\ExtensionTok{quarto}\NormalTok{ install tinytex}
\ExtensionTok{quarto}\NormalTok{ render tutorial{-}dashboard{-}emg.qmd }\AttributeTok{{-}{-}to}\NormalTok{ pdf}
\end{Highlighting}
\end{Shaded}

La primera orden solo es necesaria si falta una distribución TeX. El QMD
contiene su estilo y todo el código; los bloques se imprimen, no se
ejecutan al compilar. La composición no requiere descargar el dataset.
La aplicación sí necesita los módulos y los archivos, o conexión para
obtenerlos.

\subsection{Responsabilidad de cada
módulo}\label{responsabilidad-de-cada-muxf3dulo}

{\def\LTcaptype{none} % do not increment counter
\begin{longtable}[]{@{}
  >{\raggedright\arraybackslash}p{(\linewidth - 2\tabcolsep) * \real{0.5000}}
  >{\raggedright\arraybackslash}p{(\linewidth - 2\tabcolsep) * \real{0.5000}}@{}}
\toprule\noalign{}
\begin{minipage}[b]{\linewidth}\raggedright
Archivo
\end{minipage} & \begin{minipage}[b]{\linewidth}\raggedright
Función
\end{minipage} \\
\midrule\noalign{}
\endhead
\bottomrule\noalign{}
\endlastfoot
\texttt{datos.py} & Descargar, verificar, calibrar y describir los
registros \\
\texttt{motor.py} & Filtrar y calcular indicadores sin depender de la
interfaz \\
\texttt{app.py} & Gestionar selecciones, figuras, revisión y
exportación \\
\texttt{verificar.py} & Comprobar integridad y propiedades del
cálculo \\
\end{longtable}
}

Conserve \texttt{pyproject.toml} y \texttt{uv.lock}. Al compartir,
incluya también la atribución del dataset. Cada exportación identifica
registro, hash, parámetros e intervalo para que un tercero pueda
reproducir el resultado.

\section{Implementar el acceso al
dataset}\label{implementar-el-acceso-al-dataset}

\subsection{Fuente y hashes
publicados}\label{fuente-y-hashes-publicados}

Cree \texttt{datos.py}. Este bloque fija la versión y el conjunto de
archivos autorizado. Los hashes provienen del listado del repositorio,
no de una copia local asumida correcta.

\begin{Shaded}
\begin{Highlighting}[]
\CommentTok{"""Acceso a EMGDB 1.0.0. Lector limitado a sus archivos WFDB 16."""}
\ImportTok{from}\NormalTok{ pathlib }\ImportTok{import}\NormalTok{ Path}
\ImportTok{from}\NormalTok{ urllib.request }\ImportTok{import}\NormalTok{ urlopen}
\ImportTok{import}\NormalTok{ hashlib}
\ImportTok{import}\NormalTok{ numpy }\ImportTok{as}\NormalTok{ np}

\NormalTok{BASE }\OperatorTok{=} \StringTok{\textquotesingle{}https://physionet.org/files/emgdb/1.0.0/\textquotesingle{}}
\NormalTok{DOI }\OperatorTok{=} \StringTok{\textquotesingle{}https://doi.org/10.13026/C24S3D\textquotesingle{}}
\NormalTok{CARPETA }\OperatorTok{=}\NormalTok{ Path(}\VariableTok{\_\_file\_\_}\NormalTok{).resolve().parent }\OperatorTok{/} \StringTok{\textquotesingle{}datos\textquotesingle{}}
\NormalTok{ETIQUETAS }\OperatorTok{=}\NormalTok{ \{}
    \StringTok{\textquotesingle{}emg\_neuropathy\textquotesingle{}}\NormalTok{: }\StringTok{\textquotesingle{}Neuropatía por radiculopatía L5\textquotesingle{}}\NormalTok{,}
    \StringTok{\textquotesingle{}emg\_healthy\textquotesingle{}}\NormalTok{: }\StringTok{\textquotesingle{}Sin historia de enfermedad neuromuscular\textquotesingle{}}\NormalTok{,}
    \StringTok{\textquotesingle{}emg\_myopathy\textquotesingle{}}\NormalTok{: }\StringTok{\textquotesingle{}Miopatía por polimiositis\textquotesingle{}}\NormalTok{,}
\NormalTok{\}}
\CommentTok{\# Valores publicados en SHA256SUMS.txt de EMGDB 1.0.0.}
\NormalTok{HASHES }\OperatorTok{=}\NormalTok{ \{}
 \StringTok{\textquotesingle{}emg\_healthy.dat\textquotesingle{}}\NormalTok{:}
 \StringTok{\textquotesingle{}4990f34f0ffb314f6ce70c1d2089b7ce41fa96d19fcd9fb334815ca5033966fd\textquotesingle{}}\NormalTok{,}
 \StringTok{\textquotesingle{}emg\_healthy.hea\textquotesingle{}}\NormalTok{:}
 \StringTok{\textquotesingle{}61221ce33e167d246ec2ea8984ce2f25882b32b4d90bd7ef785c7b56ef95bac5\textquotesingle{}}\NormalTok{,}
 \StringTok{\textquotesingle{}emg\_myopathy.dat\textquotesingle{}}\NormalTok{:}
 \StringTok{\textquotesingle{}2c310a0f61843a5a5e906d413b6ff7f51370a91e516344af511c7fe6da9eaf7f\textquotesingle{}}\NormalTok{,}
 \StringTok{\textquotesingle{}emg\_myopathy.hea\textquotesingle{}}\NormalTok{:}
 \StringTok{\textquotesingle{}752093a1e3af87cd87e39f553b8725133f57a75413661208e5d93c8d2c60e73b\textquotesingle{}}\NormalTok{,}
 \StringTok{\textquotesingle{}emg\_neuropathy.dat\textquotesingle{}}\NormalTok{:}
 \StringTok{\textquotesingle{}108c6c71d3792f3dba645539bd9241119fb4c6a3a0f3e515399694b072375dd5\textquotesingle{}}\NormalTok{,}
 \StringTok{\textquotesingle{}emg\_neuropathy.hea\textquotesingle{}}\NormalTok{:}
 \StringTok{\textquotesingle{}a424c80b6a5cc5733456f0848c2650321cac3658d9b02b845e08b619b2807ce0\textquotesingle{}}\NormalTok{,}
\NormalTok{\}}
\end{Highlighting}
\end{Shaded}

\subsection{Descargar y verificar antes de
interpretar}\label{descargar-y-verificar-antes-de-interpretar}

Añada esta función. Primero comprueba si existe una copia local; en caso
contrario descarga. Solo escribe una descarga que coincide con el hash
publicado. No se omite la comprobación para hacer desaparecer un error.

\begin{Shaded}
\begin{Highlighting}[]
\KeywordTok{def}\NormalTok{ asegurar(nombre):}
    \ControlFlowTok{if}\NormalTok{ nombre }\KeywordTok{not} \KeywordTok{in}\NormalTok{ HASHES:}
        \ControlFlowTok{raise} \PreprocessorTok{ValueError}\NormalTok{(}\StringTok{\textquotesingle{}Archivo fuera del dataset permitido.\textquotesingle{}}\NormalTok{)}
\NormalTok{    ruta }\OperatorTok{=}\NormalTok{ CARPETA }\OperatorTok{/}\NormalTok{ nombre}
\NormalTok{    CARPETA.mkdir(exist\_ok}\OperatorTok{=}\VariableTok{True}\NormalTok{)}
    \ControlFlowTok{if}\NormalTok{ ruta.exists():}
\NormalTok{        contenido }\OperatorTok{=}\NormalTok{ ruta.read\_bytes()}
    \ControlFlowTok{else}\NormalTok{:}
        \ControlFlowTok{with}\NormalTok{ urlopen(BASE }\OperatorTok{+}\NormalTok{ nombre, timeout}\OperatorTok{=}\DecValTok{60}\NormalTok{) }\ImportTok{as}\NormalTok{ respuesta:}
\NormalTok{            contenido }\OperatorTok{=}\NormalTok{ respuesta.read()}
\NormalTok{    real }\OperatorTok{=}\NormalTok{ hashlib.sha256(contenido).hexdigest()}
    \ControlFlowTok{if}\NormalTok{ real }\OperatorTok{!=}\NormalTok{ HASHES[nombre]:}
        \ControlFlowTok{raise} \PreprocessorTok{ValueError}\NormalTok{(}\SpecialStringTok{f\textquotesingle{}Integridad incorrecta: }\SpecialCharTok{\{}\NormalTok{nombre}\SpecialCharTok{\}}\SpecialStringTok{\textquotesingle{}}\NormalTok{)}
    \ControlFlowTok{if} \KeywordTok{not}\NormalTok{ ruta.exists():}
\NormalTok{        ruta.write\_bytes(contenido)}
    \ControlFlowTok{return}\NormalTok{ contenido}
\end{Highlighting}
\end{Shaded}

\subsection{Leer cuentas digitales y obtener
microvoltios}\label{leer-cuentas-digitales-y-obtener-microvoltios}

Para estos archivos, WFDB 16 representa enteros con signo de 16 bits en
orden little-endian. La cabecera especifica la ganancia y el cero. La
conversión es

\[x_{\mu V}[n]=\frac{d[n]-d_0}{G}\,1000,
\qquad G=10000\ \text{cuentas/mV}.\]

Cada cuenta representa 0.1 \(\mu V\) en esta calibración. El tiempo se
reconstruye como \(t[n]=n/f_s\). Este origen en cero es una convención
de análisis; no es una fecha de adquisición. Añada:

\begin{Shaded}
\begin{Highlighting}[]
\KeywordTok{def}\NormalTok{ cargar(registro):}
    \ControlFlowTok{if}\NormalTok{ registro }\KeywordTok{not} \KeywordTok{in}\NormalTok{ ETIQUETAS:}
        \ControlFlowTok{raise} \PreprocessorTok{ValueError}\NormalTok{(}\StringTok{\textquotesingle{}Registro desconocido.\textquotesingle{}}\NormalTok{)}
\NormalTok{    texto }\OperatorTok{=}\NormalTok{ asegurar(registro }\OperatorTok{+} \StringTok{\textquotesingle{}.hea\textquotesingle{}}\NormalTok{).decode(}\StringTok{\textquotesingle{}ascii\textquotesingle{}}\NormalTok{)}
\NormalTok{    lineas }\OperatorTok{=}\NormalTok{ texto.splitlines()}
\NormalTok{    cabecera, canal }\OperatorTok{=}\NormalTok{ lineas[}\DecValTok{0}\NormalTok{].split(), lineas[}\DecValTok{1}\NormalTok{].split()}
    \ControlFlowTok{if}\NormalTok{ cabecera[}\DecValTok{1}\NormalTok{] }\OperatorTok{!=} \StringTok{\textquotesingle{}1\textquotesingle{}} \KeywordTok{or}\NormalTok{ canal[}\DecValTok{1}\NormalTok{] }\OperatorTok{!=} \StringTok{\textquotesingle{}16\textquotesingle{}}\NormalTok{:}
        \ControlFlowTok{raise} \PreprocessorTok{ValueError}\NormalTok{(}\StringTok{\textquotesingle{}Este lector requiere un canal WFDB 16.\textquotesingle{}}\NormalTok{)}
\NormalTok{    fs, n }\OperatorTok{=} \BuiltInTok{float}\NormalTok{(cabecera[}\DecValTok{2}\NormalTok{]), }\BuiltInTok{int}\NormalTok{(cabecera[}\DecValTok{3}\NormalTok{])}
    \ControlFlowTok{if}\NormalTok{ fs }\OperatorTok{!=} \DecValTok{4000} \KeywordTok{or}\NormalTok{ canal[}\DecValTok{2}\NormalTok{].lower() }\OperatorTok{!=} \StringTok{\textquotesingle{}10000/mv\textquotesingle{}}\NormalTok{:}
        \ControlFlowTok{raise} \PreprocessorTok{ValueError}\NormalTok{(}\StringTok{\textquotesingle{}Muestreo o ganancia inesperados.\textquotesingle{}}\NormalTok{)}
\NormalTok{    contenido }\OperatorTok{=}\NormalTok{ asegurar(registro }\OperatorTok{+} \StringTok{\textquotesingle{}.dat\textquotesingle{}}\NormalTok{)}
\NormalTok{    cuentas }\OperatorTok{=}\NormalTok{ np.frombuffer(contenido, dtype}\OperatorTok{=}\StringTok{\textquotesingle{}\textless{}i2\textquotesingle{}}\NormalTok{)}
    \ControlFlowTok{if} \BuiltInTok{len}\NormalTok{(cuentas) }\OperatorTok{!=}\NormalTok{ n }\KeywordTok{or}\NormalTok{ np.}\BuiltInTok{any}\NormalTok{(cuentas }\OperatorTok{==} \OperatorTok{{-}}\DecValTok{32768}\NormalTok{):}
        \ControlFlowTok{raise} \PreprocessorTok{ValueError}\NormalTok{(}\StringTok{\textquotesingle{}Longitud incorrecta o muestras faltantes.\textquotesingle{}}\NormalTok{)}
\NormalTok{    cero }\OperatorTok{=} \BuiltInTok{int}\NormalTok{(canal[}\DecValTok{4}\NormalTok{])}
\NormalTok{    x }\OperatorTok{=}\NormalTok{ (cuentas.astype(}\BuiltInTok{float}\NormalTok{) }\OperatorTok{{-}}\NormalTok{ cero) }\OperatorTok{/} \DecValTok{10000} \OperatorTok{*} \DecValTok{1000}
\NormalTok{    t }\OperatorTok{=}\NormalTok{ np.arange(n) }\OperatorTok{/}\NormalTok{ fs}
\NormalTok{    meta }\OperatorTok{=} \BuiltInTok{dict}\NormalTok{(registro}\OperatorTok{=}\NormalTok{registro, etiqueta}\OperatorTok{=}\NormalTok{ETIQUETAS[registro],}
\NormalTok{                fs\_hz}\OperatorTok{=}\NormalTok{fs, n}\OperatorTok{=}\NormalTok{n, duracion\_s}\OperatorTok{=}\NormalTok{n}\OperatorTok{/}\NormalTok{fs, unidades}\OperatorTok{=}\StringTok{\textquotesingle{}uV\textquotesingle{}}\NormalTok{,}
\NormalTok{                modalidad}\OperatorTok{=}\StringTok{\textquotesingle{}EMG intramuscular\textquotesingle{}}\NormalTok{, musculo}\OperatorTok{=}\StringTok{\textquotesingle{}Tibial anterior\textquotesingle{}}\NormalTok{,}
\NormalTok{                doi}\OperatorTok{=}\NormalTok{DOI, licencia}\OperatorTok{=}\StringTok{\textquotesingle{}ODC{-}By 1.0\textquotesingle{}}\NormalTok{,}
\NormalTok{                sha256\_dat}\OperatorTok{=}\NormalTok{HASHES[registro }\OperatorTok{+} \StringTok{\textquotesingle{}.dat\textquotesingle{}}\NormalTok{],}
\NormalTok{                sha256\_hea}\OperatorTok{=}\NormalTok{HASHES[registro }\OperatorTok{+} \StringTok{\textquotesingle{}.hea\textquotesingle{}}\NormalTok{])}
    \ControlFlowTok{return}\NormalTok{ t, x, meta}


\ControlFlowTok{if} \VariableTok{\_\_name\_\_} \OperatorTok{==} \StringTok{\textquotesingle{}\_\_main\_\_\textquotesingle{}}\NormalTok{:}
    \ControlFlowTok{for}\NormalTok{ registro }\KeywordTok{in}\NormalTok{ ETIQUETAS:}
\NormalTok{        \_, \_, meta }\OperatorTok{=}\NormalTok{ cargar(registro)}
        \BuiltInTok{print}\NormalTok{(registro, meta[}\StringTok{\textquotesingle{}n\textquotesingle{}}\NormalTok{], meta[}\StringTok{\textquotesingle{}duracion\_s\textquotesingle{}}\NormalTok{])}
\end{Highlighting}
\end{Shaded}

Ejecutar \texttt{datos.py} debe confirmar las longitudes y duraciones de
la tabla del dataset. \textbf{No} interprete las dos columnas de un
archivo de texto como dos músculos: una puede ser tiempo y la otra
señal. Esta versión evita esa ambigüedad leyendo la cabecera y el
binario.

\section{Implementar los indicadores}\label{implementar-los-indicadores}

\subsection{Validar y filtrar}\label{validar-y-filtrar}

Cree \texttt{motor.py} con este bloque. Se rechazan entradas no finitas
y canales completamente planos. Esto no detecta todos los artefactos: es
una primera condición de entrada.

\begin{Shaded}
\begin{Highlighting}[]
\CommentTok{"""Descriptores exploratorios; no clasificador de enfermedad."""}
\ImportTok{import}\NormalTok{ numpy }\ImportTok{as}\NormalTok{ np}
\ImportTok{import}\NormalTok{ pandas }\ImportTok{as}\NormalTok{ pd}
\ImportTok{from}\NormalTok{ scipy }\ImportTok{import}\NormalTok{ signal}
\ImportTok{from}\NormalTok{ scipy.integrate }\ImportTok{import}\NormalTok{ cumulative\_trapezoid}


\KeywordTok{def}\NormalTok{ preparar(x, fs):}
\NormalTok{    x }\OperatorTok{=}\NormalTok{ np.asarray(x, dtype}\OperatorTok{=}\BuiltInTok{float}\NormalTok{)}
    \ControlFlowTok{if}\NormalTok{ x.ndim }\OperatorTok{!=} \DecValTok{1} \KeywordTok{or} \KeywordTok{not}\NormalTok{ np.isfinite(x).}\BuiltInTok{all}\NormalTok{():}
        \ControlFlowTok{raise} \PreprocessorTok{ValueError}\NormalTok{(}\StringTok{\textquotesingle{}Se requiere un vector finito.\textquotesingle{}}\NormalTok{)}
    \ControlFlowTok{if}\NormalTok{ fs }\OperatorTok{!=} \DecValTok{4000} \KeywordTok{or} \BuiltInTok{len}\NormalTok{(x) }\OperatorTok{\textless{}} \DecValTok{2}\OperatorTok{*}\NormalTok{fs:}
        \ControlFlowTok{raise} \PreprocessorTok{ValueError}\NormalTok{(}\StringTok{\textquotesingle{}Se requieren 4000 Hz y al menos 2 s.\textquotesingle{}}\NormalTok{)}
    \ControlFlowTok{if}\NormalTok{ np.ptp(x) }\OperatorTok{\textless{}} \FloatTok{1e{-}9}\NormalTok{:}
        \ControlFlowTok{raise} \PreprocessorTok{ValueError}\NormalTok{(}\StringTok{\textquotesingle{}Canal plano: no interpretable.\textquotesingle{}}\NormalTok{)}
\NormalTok{    sos }\OperatorTok{=}\NormalTok{ signal.butter(}\DecValTok{4}\NormalTok{, [}\DecValTok{20}\NormalTok{, }\DecValTok{1000}\NormalTok{], fs}\OperatorTok{=}\NormalTok{fs,}
\NormalTok{                        btype}\OperatorTok{=}\StringTok{\textquotesingle{}bandpass\textquotesingle{}}\NormalTok{, output}\OperatorTok{=}\StringTok{\textquotesingle{}sos\textquotesingle{}}\NormalTok{)}
    \ControlFlowTok{return}\NormalTok{ signal.sosfiltfilt(sos, x)}
\end{Highlighting}
\end{Shaded}

\subsection{Una función para una
ventana}\label{una-funciuxf3n-para-una-ventana}

Añada el bloque siguiente. La misma función alimenta tarjetas,
tendencias y verificaciones; así se evita calcular una fórmula diferente
para cada componente del tablero.

\begin{Shaded}
\begin{Highlighting}[]
\KeywordTok{def}\NormalTok{ medir(y, fs):}
\NormalTok{    y }\OperatorTok{=}\NormalTok{ np.asarray(y, dtype}\OperatorTok{=}\BuiltInTok{float}\NormalTok{)}
    \ControlFlowTok{if}\NormalTok{ y.ndim }\OperatorTok{!=} \DecValTok{1} \KeywordTok{or} \BuiltInTok{len}\NormalTok{(y) }\OperatorTok{\textless{}} \BuiltInTok{int}\NormalTok{(}\FloatTok{0.5}\OperatorTok{*}\NormalTok{fs):}
        \ControlFlowTok{raise} \PreprocessorTok{ValueError}\NormalTok{(}\StringTok{\textquotesingle{}Ventana demasiado corta.\textquotesingle{}}\NormalTok{)}
    \ControlFlowTok{if} \KeywordTok{not}\NormalTok{ np.isfinite(y).}\BuiltInTok{all}\NormalTok{():}
        \ControlFlowTok{raise} \PreprocessorTok{ValueError}\NormalTok{(}\StringTok{\textquotesingle{}Ventana con valores no finitos.\textquotesingle{}}\NormalTok{)}
\NormalTok{    z }\OperatorTok{=}\NormalTok{ y }\OperatorTok{{-}}\NormalTok{ np.mean(y)}
\NormalTok{    rms }\OperatorTok{=} \BuiltInTok{float}\NormalTok{(np.sqrt(np.mean(z}\OperatorTok{*}\NormalTok{z)))}
\NormalTok{    f, p }\OperatorTok{=}\NormalTok{ signal.welch(z, fs}\OperatorTok{=}\NormalTok{fs, window}\OperatorTok{=}\StringTok{\textquotesingle{}hann\textquotesingle{}}\NormalTok{,}
\NormalTok{                        nperseg}\OperatorTok{=}\BuiltInTok{int}\NormalTok{(}\FloatTok{0.25}\OperatorTok{*}\NormalTok{fs),}
\NormalTok{                        noverlap}\OperatorTok{=}\BuiltInTok{int}\NormalTok{(}\FloatTok{0.125}\OperatorTok{*}\NormalTok{fs),}
\NormalTok{                        detrend}\OperatorTok{=}\StringTok{\textquotesingle{}constant\textquotesingle{}}\NormalTok{, scaling}\OperatorTok{=}\StringTok{\textquotesingle{}density\textquotesingle{}}\NormalTok{)}
\NormalTok{    m }\OperatorTok{=}\NormalTok{ (f }\OperatorTok{\textgreater{}=} \DecValTok{20}\NormalTok{) }\OperatorTok{\&}\NormalTok{ (f }\OperatorTok{\textless{}=} \DecValTok{1000}\NormalTok{)}
\NormalTok{    fb, pb }\OperatorTok{=}\NormalTok{ f[m], p[m]}
\NormalTok{    acumulada }\OperatorTok{=}\NormalTok{ cumulative\_trapezoid(pb, fb, initial}\OperatorTok{=}\DecValTok{0}\NormalTok{)}
    \ControlFlowTok{if}\NormalTok{ acumulada[}\OperatorTok{{-}}\DecValTok{1}\NormalTok{] }\OperatorTok{\textless{}=} \FloatTok{1e{-}12}\NormalTok{:}
\NormalTok{        mediana }\OperatorTok{=}\NormalTok{ np.nan}
    \ControlFlowTok{else}\NormalTok{:}
\NormalTok{        mediana }\OperatorTok{=} \BuiltInTok{float}\NormalTok{(np.interp(acumulada[}\OperatorTok{{-}}\DecValTok{1}\NormalTok{]}\OperatorTok{/}\DecValTok{2}\NormalTok{, acumulada, fb))}
\NormalTok{    metricas }\OperatorTok{=} \BuiltInTok{dict}\NormalTok{(rms\_uv}\OperatorTok{=}\NormalTok{rms, pico\_pico\_uv}\OperatorTok{=}\BuiltInTok{float}\NormalTok{(np.ptp(z)),}
\NormalTok{                   cresta}\OperatorTok{=}\BuiltInTok{float}\NormalTok{(np.}\BuiltInTok{max}\NormalTok{(np.}\BuiltInTok{abs}\NormalTok{(z))}\OperatorTok{/}\NormalTok{rms)}
                   \ControlFlowTok{if}\NormalTok{ rms }\OperatorTok{\textgreater{}} \FloatTok{1e{-}9} \ControlFlowTok{else}\NormalTok{ np.nan,}
\NormalTok{                   fmed\_20\_1000\_hz}\OperatorTok{=}\NormalTok{mediana)}
    \ControlFlowTok{return}\NormalTok{ metricas, f, p}
\end{Highlighting}
\end{Shaded}

Un \texttt{NaN} de frecuencia mediana indica potencia insuficiente para
definirla según la tolerancia numérica, no ausencia de enfermedad. No
reemplace automáticamente ese valor por cero.

\subsection{Recorrer intervalos sin perder sus
límites}\label{recorrer-intervalos-sin-perder-sus-luxedmites}

Complete el módulo con \texttt{ventanas}. Convertir a índices enteros
evita discrepancias debidas a comparaciones flotantes en los límites.
Toda fila conserva inicio, final y centro temporal.

\begin{Shaded}
\begin{Highlighting}[]
\KeywordTok{def}\NormalTok{ ventanas(y, fs, inicio, fin, W):}
    \CommentTok{\# Índices enteros: intervalos [i, j), sin ambigüedad de redondeo.}
\NormalTok{    N }\OperatorTok{=} \BuiltInTok{round}\NormalTok{(W}\OperatorTok{*}\NormalTok{fs)}
\NormalTok{    paso }\OperatorTok{=}\NormalTok{ N}\OperatorTok{//}\DecValTok{2}
\NormalTok{    i0, i1 }\OperatorTok{=} \BuiltInTok{round}\NormalTok{(inicio}\OperatorTok{*}\NormalTok{fs), }\BuiltInTok{round}\NormalTok{(fin}\OperatorTok{*}\NormalTok{fs)}
    \ControlFlowTok{if} \KeywordTok{not}\NormalTok{ (}\DecValTok{0} \OperatorTok{\textless{}=}\NormalTok{ i0 }\OperatorTok{\textless{}}\NormalTok{ i1 }\OperatorTok{\textless{}=} \BuiltInTok{len}\NormalTok{(y)) }\KeywordTok{or}\NormalTok{ i1}\OperatorTok{{-}}\NormalTok{i0 }\OperatorTok{\textless{}}\NormalTok{ N:}
        \ControlFlowTok{raise} \PreprocessorTok{ValueError}\NormalTok{(}\StringTok{\textquotesingle{}Intervalo insuficiente o fuera del registro.\textquotesingle{}}\NormalTok{)}
\NormalTok{    rows }\OperatorTok{=}\NormalTok{ []}
    \ControlFlowTok{for}\NormalTok{ i }\KeywordTok{in} \BuiltInTok{range}\NormalTok{(i0, i1}\OperatorTok{{-}}\NormalTok{N}\OperatorTok{+}\DecValTok{1}\NormalTok{, paso):}
\NormalTok{        j }\OperatorTok{=}\NormalTok{ i}\OperatorTok{+}\NormalTok{N}
\NormalTok{        valores, \_, \_ }\OperatorTok{=}\NormalTok{ medir(y[i:j], fs)}
\NormalTok{        rows.append(}\BuiltInTok{dict}\NormalTok{(i}\OperatorTok{=}\NormalTok{i, j}\OperatorTok{=}\NormalTok{j, inicio\_s}\OperatorTok{=}\NormalTok{i}\OperatorTok{/}\NormalTok{fs, fin\_s}\OperatorTok{=}\NormalTok{j}\OperatorTok{/}\NormalTok{fs,}
\NormalTok{                         centro\_s}\OperatorTok{=}\NormalTok{(i}\OperatorTok{+}\NormalTok{j)}\OperatorTok{/}\NormalTok{(}\DecValTok{2}\OperatorTok{*}\NormalTok{fs), }\OperatorTok{**}\NormalTok{valores))}
    \ControlFlowTok{return}\NormalTok{ pd.DataFrame(rows)}
\end{Highlighting}
\end{Shaded}

\textbf{Actividad breve.} Use el mismo intervalo con ventanas de 0.5 y 2
s. Explique por qué cambian la cantidad de puntos y su estabilidad. No
seleccione la ventana solo porque produce una curva visualmente más
suave.

\section{Construir el dashboard}\label{construir-el-dashboard}

\subsection{Diseñar para la lectura
clínica}\label{diseuxf1ar-para-la-lectura-cluxednica}

La interfaz muestra primero fuente y etiqueta, después la señal y
finalmente los indicadores. El médico debe poder inspeccionar la
evidencia que produjo cada cifra. Los colores distinguen original y
filtrada, sin codificar severidad.

\Needspace{12\baselineskip}

{\def\LTcaptype{none} % do not increment counter
\begin{longtable}[]{@{}
  >{\raggedright\arraybackslash}p{(\linewidth - 2\tabcolsep) * \real{0.5000}}
  >{\raggedright\arraybackslash}p{(\linewidth - 2\tabcolsep) * \real{0.5000}}@{}}
\toprule\noalign{}
\begin{minipage}[b]{\linewidth}\raggedright
Componente
\end{minipage} & \begin{minipage}[b]{\linewidth}\raggedright
Pregunta que resuelve
\end{minipage} \\
\midrule\noalign{}
\endhead
\bottomrule\noalign{}
\endlastfoot
Selector de registro & ¿Qué ejemplo clínico estoy observando? \\
Intervalo y ventana & ¿Qué muestras entran en el cálculo? \\
Original y filtrada & ¿El procesamiento conserva lo que quiero
estudiar? \\
Tarjetas & ¿Qué describen amplitud, extremos y espectro? \\
Detalle y PSD & ¿Qué evidencia explica las tarjetas? \\
Tendencia & ¿Cómo varía el descriptor dentro del registro? \\
Exportación & ¿Puedo repetir exactamente este análisis? \\
\end{longtable}
}

El prototipo muestra todas las muestras del intervalo: evita que una
reducción visual arbitraria oculte extremos. Para archivos muy grandes
sería necesario otro diseño de visualización; reducir puntos en pantalla
nunca debe alterar los datos usados en las métricas.

\subsection{Cargar y mostrar el
contexto}\label{cargar-y-mostrar-el-contexto}

Cree \texttt{app.py} con el siguiente bloque. La selección inicial es el
registro de neuropatía. La etiqueta se toma de metadatos; no se presenta
como salida de un clasificador.

\begin{Shaded}
\begin{Highlighting}[]
\ImportTok{import}\NormalTok{ json}
\ImportTok{import}\NormalTok{ plotly.graph\_objects }\ImportTok{as}\NormalTok{ go}
\ImportTok{import}\NormalTok{ streamlit }\ImportTok{as}\NormalTok{ st}
\ImportTok{from}\NormalTok{ datos }\ImportTok{import}\NormalTok{ cargar, ETIQUETAS}
\ImportTok{from}\NormalTok{ motor }\ImportTok{import}\NormalTok{ preparar, medir, ventanas}

\NormalTok{st.set\_page\_config(page\_title}\OperatorTok{=}\StringTok{\textquotesingle{}EMGDB | Dashboard\textquotesingle{}}\NormalTok{, layout}\OperatorTok{=}\StringTok{\textquotesingle{}wide\textquotesingle{}}\NormalTok{)}
\NormalTok{st.title(}\StringTok{\textquotesingle{}Dashboard de EMG y radiculopatía L5\textquotesingle{}}\NormalTok{)}
\NormalTok{st.caption(}\StringTok{\textquotesingle{}Dataset abierto EMGDB 1.0.0 | Análisis retrospectivo docente\textquotesingle{}}\NormalTok{)}
\NormalTok{st.info(}\StringTok{\textquotesingle{}Tres personas distintas. No son visitas de un mismo paciente.\textquotesingle{}}\NormalTok{)}
\NormalTok{registro }\OperatorTok{=}\NormalTok{ st.sidebar.selectbox(}\StringTok{\textquotesingle{}Registro\textquotesingle{}}\NormalTok{, }\BuiltInTok{list}\NormalTok{(ETIQUETAS),}
\NormalTok{                               format\_func}\OperatorTok{=}\NormalTok{ETIQUETAS.get)}
\ControlFlowTok{try}\NormalTok{:}
\NormalTok{    t, x, meta }\OperatorTok{=}\NormalTok{ cargar(registro)}
\NormalTok{    fs }\OperatorTok{=}\NormalTok{ meta[}\StringTok{\textquotesingle{}fs\_hz\textquotesingle{}}\NormalTok{]}
\NormalTok{    y }\OperatorTok{=}\NormalTok{ preparar(x, fs)}
\ControlFlowTok{except} \PreprocessorTok{Exception} \ImportTok{as}\NormalTok{ exc:}
\NormalTok{    st.error(}\SpecialStringTok{f\textquotesingle{}No se pudo cargar o procesar: }\SpecialCharTok{\{}\NormalTok{exc}\SpecialCharTok{\}}\SpecialStringTok{\textquotesingle{}}\NormalTok{)}
\NormalTok{    st.stop()}
\NormalTok{st.caption(}\SpecialStringTok{f"}\SpecialCharTok{\{}\NormalTok{registro}\SpecialCharTok{\}}\SpecialStringTok{ | }\SpecialCharTok{\{}\NormalTok{meta[}\StringTok{\textquotesingle{}modalidad\textquotesingle{}}\NormalTok{]}\SpecialCharTok{\}}\SpecialStringTok{ | "}
           \SpecialStringTok{f"}\SpecialCharTok{\{}\NormalTok{meta[}\StringTok{\textquotesingle{}musculo\textquotesingle{}}\NormalTok{]}\SpecialCharTok{\}}\SpecialStringTok{ | }\SpecialCharTok{\{}\NormalTok{fs}\SpecialCharTok{:.0f\}}\SpecialStringTok{ Hz | }\SpecialCharTok{\{}\BuiltInTok{len}\NormalTok{(x)}\OperatorTok{/}\NormalTok{fs}\SpecialCharTok{:.4f\}}\SpecialStringTok{ s"}\NormalTok{)}
\NormalTok{st.markdown(}\StringTok{\textquotesingle{}[Fuente y licencia](https://physionet.org/content/emgdb/1.0.0/)\textquotesingle{}}\NormalTok{)}
\end{Highlighting}
\end{Shaded}

\subsection{Elegir intervalo y revisar
calidad}\label{elegir-intervalo-y-revisar-calidad}

Añada este bloque. La confirmación visual se reinicia al cambiar el
registro o los límites. La anotación describe el segmento y las dudas de
calidad. Es responsabilidad del usuario no atribuir significado
fisiológico a artefactos; el programa no incluye un detector validado de
ellos.

\begin{Shaded}
\begin{Highlighting}[]
\CommentTok{\# Descarte de 0.25 s en cada extremo del registro ya filtrado.}
\NormalTok{a, b }\OperatorTok{=}\NormalTok{ st.sidebar.slider(}\StringTok{\textquotesingle{}Intervalo de análisis (s)\textquotesingle{}}\NormalTok{,}
    \FloatTok{0.25}\NormalTok{, }\BuiltInTok{float}\NormalTok{(}\BuiltInTok{len}\NormalTok{(x)}\OperatorTok{/}\NormalTok{fs}\OperatorTok{{-}}\FloatTok{0.25}\NormalTok{), (}\FloatTok{0.25}\NormalTok{, }\FloatTok{5.25}\NormalTok{), step}\OperatorTok{=}\FloatTok{0.25}\NormalTok{,}
\NormalTok{    key}\OperatorTok{=}\StringTok{\textquotesingle{}intervalo\_\textquotesingle{}} \OperatorTok{+}\NormalTok{ registro)}
\NormalTok{W }\OperatorTok{=}\NormalTok{ st.sidebar.selectbox(}\StringTok{\textquotesingle{}Duración de ventana (s)\textquotesingle{}}\NormalTok{, [}\FloatTok{0.5}\NormalTok{, }\FloatTok{1.0}\NormalTok{, }\FloatTok{2.0}\NormalTok{], index}\OperatorTok{=}\DecValTok{1}\NormalTok{)}
\NormalTok{st.sidebar.caption(}\StringTok{\textquotesingle{}Paso = W/2. Banda 20–1000 Hz. Welch 0.25 s.\textquotesingle{}}\NormalTok{)}
\NormalTok{nota }\OperatorTok{=}\NormalTok{ st.sidebar.text\_area(}\StringTok{\textquotesingle{}Anotación de calidad\textquotesingle{}}\NormalTok{,}
\NormalTok{                           key}\OperatorTok{=}\StringTok{\textquotesingle{}nota\_\textquotesingle{}} \OperatorTok{+}\NormalTok{ registro)}
\NormalTok{revisado }\OperatorTok{=}\NormalTok{ st.sidebar.checkbox(}\StringTok{\textquotesingle{}Inspeccioné el intervalo seleccionado\textquotesingle{}}\NormalTok{,}
\NormalTok{                              key}\OperatorTok{=}\SpecialStringTok{f\textquotesingle{}revision\_}\SpecialCharTok{\{}\NormalTok{registro}\SpecialCharTok{\}}\SpecialStringTok{\_}\SpecialCharTok{\{}\NormalTok{a}\SpecialCharTok{\}}\SpecialStringTok{\_}\SpecialCharTok{\{}\NormalTok{b}\SpecialCharTok{\}}\SpecialStringTok{\textquotesingle{}}\NormalTok{)}
\ControlFlowTok{if}\NormalTok{ b}\OperatorTok{{-}}\NormalTok{a }\OperatorTok{\textless{}}\NormalTok{ W:}
\NormalTok{    st.error(}\StringTok{\textquotesingle{}Seleccione un intervalo al menos igual a la ventana.\textquotesingle{}}\NormalTok{)}
\NormalTok{    st.stop()}

\NormalTok{fig }\OperatorTok{=}\NormalTok{ go.Figure()}
\NormalTok{i0, i1 }\OperatorTok{=} \BuiltInTok{round}\NormalTok{(a}\OperatorTok{*}\NormalTok{fs), }\BuiltInTok{round}\NormalTok{(b}\OperatorTok{*}\NormalTok{fs)}
\NormalTok{fig.add\_trace(go.Scattergl(x}\OperatorTok{=}\NormalTok{t[i0:i1], y}\OperatorTok{=}\NormalTok{x[i0:i1], name}\OperatorTok{=}\StringTok{\textquotesingle{}Original\textquotesingle{}}\NormalTok{))}
\NormalTok{fig.add\_trace(go.Scattergl(x}\OperatorTok{=}\NormalTok{t[i0:i1], y}\OperatorTok{=}\NormalTok{y[i0:i1], name}\OperatorTok{=}\StringTok{\textquotesingle{}Filtrada\textquotesingle{}}\NormalTok{))}
\NormalTok{fig.update\_layout(title}\OperatorTok{=}\StringTok{\textquotesingle{}Control de calidad: todas las muestras del intervalo\textquotesingle{}}\NormalTok{,}
\NormalTok{                  xaxis\_title}\OperatorTok{=}\StringTok{\textquotesingle{}Tiempo (s)\textquotesingle{}}\NormalTok{, yaxis\_title}\OperatorTok{=}\StringTok{\textquotesingle{}EMG (µV)\textquotesingle{}}\NormalTok{)}
\NormalTok{st.plotly\_chart(fig, width}\OperatorTok{=}\StringTok{\textquotesingle{}stretch\textquotesingle{}}\NormalTok{)}
\NormalTok{st.caption(}\StringTok{\textquotesingle{}Use zoom para inspeccionar eventos. La revisión visual \textquotesingle{}}
           \StringTok{\textquotesingle{}no sustituye anotaciones expertas ni detecta todo artefacto.\textquotesingle{}}\NormalTok{)}
\ControlFlowTok{if} \KeywordTok{not}\NormalTok{ revisado:}
\NormalTok{    st.warning(}\StringTok{\textquotesingle{}Revise original y filtrada antes de habilitar indicadores.\textquotesingle{}}\NormalTok{)}
\NormalTok{    st.stop()}
\end{Highlighting}
\end{Shaded}

\subsection{Vincular tarjetas y ventana
activa}\label{vincular-tarjetas-y-ventana-activa}

El índice seleccionado determina las cuatro tarjetas, el detalle
temporal y la PSD. Cada cifra corresponde al mismo intervalo. El pico a
pico se rotula explícitamente como medida de ventana.

\begin{Shaded}
\begin{Highlighting}[]
\NormalTok{res }\OperatorTok{=}\NormalTok{ ventanas(y, fs, a, b, W)}
\NormalTok{seleccion }\OperatorTok{=}\NormalTok{ st.select\_slider(}\StringTok{\textquotesingle{}Ventana activa\textquotesingle{}}\NormalTok{, options}\OperatorTok{=}\BuiltInTok{list}\NormalTok{(res.index),}
\NormalTok{    format\_func}\OperatorTok{=}\KeywordTok{lambda}\NormalTok{ k: }\SpecialStringTok{f"[}\SpecialCharTok{\{}\NormalTok{res}\SpecialCharTok{.}\NormalTok{loc[k,}\StringTok{\textquotesingle{}inicio\_s\textquotesingle{}}\NormalTok{]}\SpecialCharTok{:.3f\}}\SpecialStringTok{, "}
                         \SpecialStringTok{f"}\SpecialCharTok{\{}\NormalTok{res}\SpecialCharTok{.}\NormalTok{loc[k,}\StringTok{\textquotesingle{}fin\_s\textquotesingle{}}\NormalTok{]}\SpecialCharTok{:.3f\}}\SpecialStringTok{) s"}\NormalTok{)}
\NormalTok{r }\OperatorTok{=}\NormalTok{ res.loc[seleccion]}
\NormalTok{i, j }\OperatorTok{=} \BuiltInTok{int}\NormalTok{(r.i), }\BuiltInTok{int}\NormalTok{(r.j)}
\NormalTok{valores, f, p }\OperatorTok{=}\NormalTok{ medir(y[i:j], fs)}
\ControlFlowTok{for}\NormalTok{ col, clave, titulo }\KeywordTok{in} \BuiltInTok{zip}\NormalTok{(st.columns(}\DecValTok{4}\NormalTok{), valores,}
\NormalTok{    [}\StringTok{\textquotesingle{}RMS (µV)\textquotesingle{}}\NormalTok{, }\StringTok{\textquotesingle{}Pico a pico de ventana (µV)\textquotesingle{}}\NormalTok{,}
     \StringTok{\textquotesingle{}Factor de cresta\textquotesingle{}}\NormalTok{, }\StringTok{\textquotesingle{}Frecuencia mediana 20–1000 Hz\textquotesingle{}}\NormalTok{]):}
\NormalTok{    col.metric(titulo, }\SpecialStringTok{f\textquotesingle{}}\SpecialCharTok{\{}\NormalTok{valores[clave]}\SpecialCharTok{:.2f\}}\SpecialStringTok{\textquotesingle{}}\NormalTok{)}
\NormalTok{st.caption(}\StringTok{\textquotesingle{}Descriptores de la señal; no puntajes de gravedad ni medidas \textquotesingle{}}
           \StringTok{\textquotesingle{}automáticas de una unidad motora aislada.\textquotesingle{}}\NormalTok{)}

\NormalTok{fig }\OperatorTok{=}\NormalTok{ go.Figure(go.Scatter(x}\OperatorTok{=}\NormalTok{t[i:j], y}\OperatorTok{=}\NormalTok{y[i:j], name}\OperatorTok{=}\StringTok{\textquotesingle{}Ventana activa\textquotesingle{}}\NormalTok{))}
\NormalTok{fig.update\_layout(title}\OperatorTok{=}\StringTok{\textquotesingle{}Detalle temporal sin reducción de muestras\textquotesingle{}}\NormalTok{,}
\NormalTok{                  xaxis\_title}\OperatorTok{=}\StringTok{\textquotesingle{}Tiempo (s)\textquotesingle{}}\NormalTok{, yaxis\_title}\OperatorTok{=}\StringTok{\textquotesingle{}EMG filtrada (µV)\textquotesingle{}}\NormalTok{)}
\NormalTok{st.plotly\_chart(fig, width}\OperatorTok{=}\StringTok{\textquotesingle{}stretch\textquotesingle{}}\NormalTok{)}
\NormalTok{fig }\OperatorTok{=}\NormalTok{ go.Figure(go.Scatter(x}\OperatorTok{=}\NormalTok{f, y}\OperatorTok{=}\NormalTok{p))}
\NormalTok{fig.update\_layout(title}\OperatorTok{=}\StringTok{\textquotesingle{}PSD de la EMG filtrada y centrada\textquotesingle{}}\NormalTok{,}
\NormalTok{                  xaxis\_title}\OperatorTok{=}\StringTok{\textquotesingle{}Frecuencia (Hz)\textquotesingle{}}\NormalTok{, yaxis\_title}\OperatorTok{=}\StringTok{\textquotesingle{}PSD (µV²/Hz)\textquotesingle{}}\NormalTok{,}
\NormalTok{                  xaxis\_range}\OperatorTok{=}\NormalTok{[}\DecValTok{0}\NormalTok{, }\DecValTok{1200}\NormalTok{])}
\NormalTok{st.plotly\_chart(fig, width}\OperatorTok{=}\StringTok{\textquotesingle{}stretch\textquotesingle{}}\NormalTok{)}
\end{Highlighting}
\end{Shaded}

\subsection{Tendencias y exportación}\label{tendencias-y-exportaciuxf3n}

Complete \texttt{app.py}. El CSV guarda los indicadores y el
identificador del registro. El JSON conserva la configuración y la
anotación. Ambos deben mantenerse juntos.

\begin{Shaded}
\begin{Highlighting}[]
\NormalTok{variable }\OperatorTok{=}\NormalTok{ st.selectbox(}\StringTok{\textquotesingle{}Indicador de tendencia\textquotesingle{}}\NormalTok{, }\BuiltInTok{list}\NormalTok{(valores))}
\NormalTok{fig }\OperatorTok{=}\NormalTok{ go.Figure(go.Scatter(x}\OperatorTok{=}\NormalTok{res.centro\_s, y}\OperatorTok{=}\NormalTok{res[variable],}
\NormalTok{                          mode}\OperatorTok{=}\StringTok{\textquotesingle{}lines+markers\textquotesingle{}}\NormalTok{))}
\NormalTok{fig.update\_layout(title}\OperatorTok{=}\StringTok{\textquotesingle{}Evolución dentro del registro\textquotesingle{}}\NormalTok{,}
\NormalTok{                  xaxis\_title}\OperatorTok{=}\StringTok{\textquotesingle{}Centro de ventana (s)\textquotesingle{}}\NormalTok{, yaxis\_title}\OperatorTok{=}\NormalTok{variable)}
\NormalTok{st.plotly\_chart(fig, width}\OperatorTok{=}\StringTok{\textquotesingle{}stretch\textquotesingle{}}\NormalTok{)}
\NormalTok{st.caption(}\StringTok{\textquotesingle{}Solapamiento 50 \%. Los puntos no son réplicas independientes.\textquotesingle{}}\NormalTok{)}

\NormalTok{meta.update(banda\_hz}\OperatorTok{=}\NormalTok{[}\DecValTok{20}\NormalTok{,}\DecValTok{1000}\NormalTok{], filtro}\OperatorTok{=}\StringTok{\textquotesingle{}Butterworth SOS bidireccional\textquotesingle{}}\NormalTok{,}
\NormalTok{            orden\_diseno}\OperatorTok{=}\DecValTok{4}\NormalTok{, borde\_s}\OperatorTok{=}\FloatTok{0.25}\NormalTok{, ventana\_s}\OperatorTok{=}\NormalTok{W, paso\_s}\OperatorTok{=}\NormalTok{W}\OperatorTok{/}\DecValTok{2}\NormalTok{,}
\NormalTok{            welch\_s}\OperatorTok{=}\FloatTok{0.25}\NormalTok{, solapamiento\_welch\_s}\OperatorTok{=}\FloatTok{0.125}\NormalTok{,}
\NormalTok{            intervalo\_s}\OperatorTok{=}\NormalTok{[i0}\OperatorTok{/}\NormalTok{fs,i1}\OperatorTok{/}\NormalTok{fs], nota\_calidad}\OperatorTok{=}\NormalTok{nota,}
\NormalTok{            inspeccion\_visual}\OperatorTok{=}\VariableTok{True}\NormalTok{, version\_pipeline}\OperatorTok{=}\StringTok{\textquotesingle{}2.0\textquotesingle{}}\NormalTok{)}
\NormalTok{res[}\StringTok{\textquotesingle{}registro\textquotesingle{}}\NormalTok{] }\OperatorTok{=}\NormalTok{ registro}
\NormalTok{res[}\StringTok{\textquotesingle{}sha256\_dat\textquotesingle{}}\NormalTok{] }\OperatorTok{=}\NormalTok{ meta[}\StringTok{\textquotesingle{}sha256\_dat\textquotesingle{}}\NormalTok{]}
\NormalTok{st.download\_button(}\StringTok{\textquotesingle{}Indicadores CSV\textquotesingle{}}\NormalTok{, res.to\_csv(index}\OperatorTok{=}\VariableTok{False}\NormalTok{).encode(),}
\NormalTok{                   registro }\OperatorTok{+} \StringTok{\textquotesingle{}\_indicadores.csv\textquotesingle{}}\NormalTok{, }\StringTok{\textquotesingle{}text/csv\textquotesingle{}}\NormalTok{)}
\NormalTok{st.download\_button(}\StringTok{\textquotesingle{}Parámetros JSON\textquotesingle{}}\NormalTok{,}
\NormalTok{                   json.dumps(meta, ensure\_ascii}\OperatorTok{=}\VariableTok{False}\NormalTok{, indent}\OperatorTok{=}\DecValTok{2}\NormalTok{),}
\NormalTok{                   registro }\OperatorTok{+} \StringTok{\textquotesingle{}\_parametros.json\textquotesingle{}}\NormalTok{, }\StringTok{\textquotesingle{}application/json\textquotesingle{}}\NormalTok{)}
\end{Highlighting}
\end{Shaded}

Streamlit vuelve a ejecutar el script cuando cambia un control. Para
estos registros pequeños, releer y verificar es aceptable. Una futura
caché debe usar el contenido o hash junto con los parámetros; no debe
reutilizar resultados de otro archivo {[}8{]}.

\section{Comprobar el resultado con datos
reales}\label{comprobar-el-resultado-con-datos-reales}

\subsection{Verificación
automática}\label{verificaciuxf3n-automuxe1tica}

Guarde este bloque como \texttt{verificar.py}. Todas las comprobaciones
parten de los archivos clínicos abiertos. La transformación de amplitud
es una prueba algebraica sobre el mismo segmento, no una nueva
observación ni un paciente artificial.

\begin{Shaded}
\begin{Highlighting}[]
\ImportTok{import}\NormalTok{ numpy }\ImportTok{as}\NormalTok{ np}
\ImportTok{from}\NormalTok{ datos }\ImportTok{import}\NormalTok{ cargar, ETIQUETAS}
\ImportTok{from}\NormalTok{ motor }\ImportTok{import}\NormalTok{ preparar, medir, ventanas}

\NormalTok{esperadas }\OperatorTok{=} \BuiltInTok{dict}\NormalTok{(emg\_healthy}\OperatorTok{=}\DecValTok{50860}\NormalTok{, emg\_myopathy}\OperatorTok{=}\DecValTok{110337}\NormalTok{,}
\NormalTok{                 emg\_neuropathy}\OperatorTok{=}\DecValTok{147858}\NormalTok{)}
\ControlFlowTok{for}\NormalTok{ registro }\KeywordTok{in}\NormalTok{ ETIQUETAS:}
\NormalTok{    t, x, meta }\OperatorTok{=}\NormalTok{ cargar(registro)  }\CommentTok{\# Incluye SHA{-}256 de DAT y HEA.}
\NormalTok{    fs }\OperatorTok{=}\NormalTok{ meta[}\StringTok{\textquotesingle{}fs\_hz\textquotesingle{}}\NormalTok{]}
    \ControlFlowTok{assert} \BuiltInTok{len}\NormalTok{(x) }\OperatorTok{==}\NormalTok{ esperadas[registro]}
    \ControlFlowTok{assert}\NormalTok{ np.allclose(np.diff(t), }\DecValTok{1}\OperatorTok{/}\NormalTok{fs)}
\NormalTok{    y }\OperatorTok{=}\NormalTok{ preparar(x, fs)}
\NormalTok{    tabla }\OperatorTok{=}\NormalTok{ ventanas(y, fs, }\DecValTok{1}\NormalTok{, }\DecValTok{3}\NormalTok{, }\DecValTok{1}\NormalTok{)}
    \ControlFlowTok{assert} \BuiltInTok{len}\NormalTok{(tabla) }\OperatorTok{==} \DecValTok{3}
\NormalTok{    m, \_, \_ }\OperatorTok{=}\NormalTok{ medir(y[}\DecValTok{4000}\NormalTok{:}\DecValTok{8000}\NormalTok{], fs)}
\NormalTok{    doble, \_, \_ }\OperatorTok{=}\NormalTok{ medir(}\DecValTok{2}\OperatorTok{*}\NormalTok{y[}\DecValTok{4000}\NormalTok{:}\DecValTok{8000}\NormalTok{], fs)}
    \ControlFlowTok{assert}\NormalTok{ np.isclose(doble[}\StringTok{\textquotesingle{}rms\_uv\textquotesingle{}}\NormalTok{], }\DecValTok{2}\OperatorTok{*}\NormalTok{m[}\StringTok{\textquotesingle{}rms\_uv\textquotesingle{}}\NormalTok{])}
    \ControlFlowTok{assert}\NormalTok{ np.isclose(doble[}\StringTok{\textquotesingle{}pico\_pico\_uv\textquotesingle{}}\NormalTok{], }\DecValTok{2}\OperatorTok{*}\NormalTok{m[}\StringTok{\textquotesingle{}pico\_pico\_uv\textquotesingle{}}\NormalTok{])}
    \ControlFlowTok{assert}\NormalTok{ np.isclose(doble[}\StringTok{\textquotesingle{}cresta\textquotesingle{}}\NormalTok{], m[}\StringTok{\textquotesingle{}cresta\textquotesingle{}}\NormalTok{])}
    \ControlFlowTok{assert}\NormalTok{ np.isclose(doble[}\StringTok{\textquotesingle{}fmed\_20\_1000\_hz\textquotesingle{}}\NormalTok{], m[}\StringTok{\textquotesingle{}fmed\_20\_1000\_hz\textquotesingle{}}\NormalTok{])}
    \BuiltInTok{print}\NormalTok{(registro, }\StringTok{\textquotesingle{}OK\textquotesingle{}}\NormalTok{, meta[}\StringTok{\textquotesingle{}n\textquotesingle{}}\NormalTok{], meta[}\StringTok{\textquotesingle{}duracion\_s\textquotesingle{}}\NormalTok{])}
\CommentTok{\# Calibración independiente de la primera muestra del registro L5.}
\NormalTok{\_, x, \_ }\OperatorTok{=}\NormalTok{ cargar(}\StringTok{\textquotesingle{}emg\_neuropathy\textquotesingle{}}\NormalTok{)}
\ControlFlowTok{assert}\NormalTok{ np.isclose(x[}\DecValTok{0}\NormalTok{], }\FloatTok{90.0}\NormalTok{)}
\NormalTok{x[}\DecValTok{10}\NormalTok{] }\OperatorTok{=}\NormalTok{ np.nan}
\ControlFlowTok{try}\NormalTok{:}
\NormalTok{    preparar(x, }\DecValTok{4000}\NormalTok{)}
\ControlFlowTok{except} \PreprocessorTok{ValueError}\NormalTok{:}
    \ControlFlowTok{pass}
\ControlFlowTok{else}\NormalTok{:}
    \ControlFlowTok{raise} \PreprocessorTok{AssertionError}\NormalTok{(}\StringTok{\textquotesingle{}Se aceptó un valor no finito.\textquotesingle{}}\NormalTok{)}
\BuiltInTok{print}\NormalTok{(}\StringTok{\textquotesingle{}Verificaciones superadas con los tres registros reales.\textquotesingle{}}\NormalTok{)}
\end{Highlighting}
\end{Shaded}

La prueba de 90 \(\mu V\) corresponde a la primera cuenta calibrada del
archivo de neuropatía. Sirve para detectar errores de unidades. El éxito
del script demuestra las propiedades verificadas, no validez clínica ni
ausencia de todos los errores posibles.

\subsection{Recorrido del usuario}\label{recorrido-del-usuario}

\begin{enumerate}
\def\labelenumi{\arabic{enumi}.}
\tightlist
\item
  Ejecute la descarga o utilice los originales del paquete. Confirme que
  los hashes pasan.
\item
  Abra el dashboard y conserve \textbf{Neuropatía por radiculopatía L5}
  como registro activo.
\item
  Use el intervalo inicial. Acerque el gráfico y compare original y
  filtrada. Anote una observación de calidad.
\item
  Habilite indicadores. Desplace la ventana activa y confirme que
  detalle, PSD y tarjetas corresponden a sus límites.
\item
  Cambie la tendencia entre RMS, pico a pico, cresta y frecuencia
  mediana. Explique por qué no representan lo mismo.
\item
  Repita con W de 0.5 y 2 s. Mantenga el intervalo para aislar el efecto
  de la ventana.
\item
  Descargue CSV y JSON. Recalcule una fila usando \texttt{medir} y sus
  índices \texttt{i}, \texttt{j}.
\item
  Abra los otros registros para contrastar ejemplos. Identifique al
  menos dos razones por las que una diferencia entre ellos no demuestra
  progresión clínica.
\end{enumerate}

\subsection{Pruebas negativas y criterios de
aceptación}\label{pruebas-negativas-y-criterios-de-aceptaciuxf3n}

{\def\LTcaptype{none} % do not increment counter
\begin{longtable}[]{@{}
  >{\raggedright\arraybackslash}p{(\linewidth - 2\tabcolsep) * \real{0.5000}}
  >{\raggedright\arraybackslash}p{(\linewidth - 2\tabcolsep) * \real{0.5000}}@{}}
\toprule\noalign{}
\begin{minipage}[b]{\linewidth}\raggedright
Perturbación
\end{minipage} & \begin{minipage}[b]{\linewidth}\raggedright
Respuesta exigida
\end{minipage} \\
\midrule\noalign{}
\endhead
\bottomrule\noalign{}
\endlastfoot
Binario cambiado o incompleto & Bloqueo por SHA-256 \\
Cabecera o ganancia inesperada & Error antes de interpretar unidades \\
Dato no finito en procesamiento & Rechazo \\
Canal completamente plano & Sin indicadores interpretables \\
Intervalo más corto que W & Mensaje para corregir selección \\
Multiplicar amplitud por dos & RMS y pico a pico duplicados; cresta y
mediana invariantes \\
Cambio de registro & Etiqueta, señal y contexto actualizados; revisión
pendiente \\
\end{longtable}
}

No se considera fallo que los tres registros no se ordenen por una
supuesta escala de gravedad. No existe esa escala en el diseño. Ajustar
filtros o ventanas hasta obtener un orden deseado sería un sesgo de
análisis.

\section{Interpretación, extensión y
evaluación}\label{interpretaciuxf3n-extensiuxf3n-y-evaluaciuxf3n}

\subsection{De la tendencia a la
afirmación}\label{de-la-tendencia-a-la-afirmaciuxf3n}

Una conclusión defendible describe \textbf{archivo, intervalo,
parámetros y resultado calculado}. Por ejemplo: ``El RMS varía entre
ventanas del registro seleccionado; el patrón debe revisarse junto con
la señal y la anotación de calidad''. Añada cifras solo después de
ejecutar el análisis.

No escriba ``la neuropatía empeora'' porque el RMS aumente, ni ``hay
menos unidades motoras'' porque crezca el factor de cresta. Las métricas
no identifican de forma única el mecanismo fisiológico. El esfuerzo, la
geometría del registro, la interferencia y el procesamiento pueden
producir cambios similares.

La comparación de tres personas sirve para explorar ejemplos. Dividir
cada registro en muchas ventanas no aumenta el número de pacientes.
Entrenar y evaluar un clasificador repartiendo al azar ventanas de estas
mismas personas induciría dependencia y no justificaría generalización
clínica.

\subsection{Qué falta para monitoreo
longitudinal}\label{quuxe9-falta-para-monitoreo-longitudinal}

El requerimiento médico de comparar visitas necesita datos nuevos con
identificadores seudónimos estables, fecha o número de visita, músculo,
lado, modalidad, calibración, tarea, nivel de activación y evaluación
clínica de referencia. Este dataset no proporciona esa estructura
longitudinal.

Como extensión de diseño, proponga una tabla \texttt{persona},
\texttt{visita}, \texttt{registro}, \texttt{musculo}, \texttt{lado},
\texttt{tarea}, \texttt{calibracion}, \texttt{configuracion},
\texttt{evaluacion\_clinica}. Programe una verificación de
compatibilidad antes de comparar. No rellene visitas inventadas con
ventanas del mismo archivo.

El tutorial analiza registros ya obtenidos. La adquisición intramuscular
corresponde al personal clínico especializado y no forma parte de la
práctica de programación.

\subsection{Preguntas para la
exposición}\label{preguntas-para-la-exposiciuxf3n}

\begin{enumerate}
\def\labelenumi{\arabic{enumi}.}
\tightlist
\item
  ¿Qué cambia al pasar de EMG superficial a intramuscular?
\item
  ¿Cómo se convierte una cuenta digital en microvoltios y dónde se
  verifica la ganancia?
\item
  ¿Cómo conecta la radiculopatía con cambios potenciales de unidades
  motoras, sin convertir un descriptor en diagnóstico?
\item
  ¿Por qué el pico a pico de ventana no es la amplitud de un potencial
  aislado?
\item
  ¿Qué diferencia hay entre ventana analítica, paso y segmento Welch?
\item
  ¿Qué información se pierde con el filtro de 1000 Hz? ¿Por qué se
  conserva el original?
\item
  ¿Por qué una frecuencia mediana diferente no demuestra fatiga o
  gravedad?
\item
  ¿Por qué las ventanas de tres personas no permiten validar seguimiento
  longitudinal?
\end{enumerate}

\subsection{Rúbrica sugerida}\label{ruxfabrica-sugerida}

Seis criterios, con nivel de 0 a 3: \textbf{3}, correcto y justificado;
\textbf{2}, correcto pero parcialmente justificado; \textbf{1},
incompleto o con errores relevantes; \textbf{0}, ausente o incorrecto.
El puntaje es \(\sum_i w_iL_i/3\) sobre 100.

\Needspace{16\baselineskip}

{\def\LTcaptype{none} % do not increment counter
\begin{longtable}[]{@{}
  >{\raggedright\arraybackslash}p{(\linewidth - 4\tabcolsep) * \real{0.3333}}
  >{\raggedleft\arraybackslash}p{(\linewidth - 4\tabcolsep) * \real{0.3333}}
  >{\raggedright\arraybackslash}p{(\linewidth - 4\tabcolsep) * \real{0.3333}}@{}}
\toprule\noalign{}
\begin{minipage}[b]{\linewidth}\raggedright
Criterio
\end{minipage} & \begin{minipage}[b]{\linewidth}\raggedleft
Peso
\end{minipage} & \begin{minipage}[b]{\linewidth}\raggedright
Evidencia exigida para nivel 3
\end{minipage} \\
\midrule\noalign{}
\endhead
\bottomrule\noalign{}
\endlastfoot
Fisiología e interpretación & 20 & Conecta lesión, señal y descriptor,
explicando los límites \\
Dataset y trazabilidad & 15 & Licencia, versión, hashes y unidades
verificadas \\
Procesamiento & 20 & Justifica banda, filtrado y escalas temporales \\
Indicadores y verificación & 20 & Fórmulas coherentes y pruebas sobre
datos reales \\
Visualización & 15 & Selecciones consistentes y evidencia asociada a
cada cifra \\
Reproducibilidad y comunicación & 10 & Ejecución documentada,
exportación y exposición clara \\
\end{longtable}
}

\section{Referencias y atribución}\label{referencias-y-atribuciuxf3n}

{[}1{]} PhysioNet. \emph{Examples of Electromyograms}, versión 1.0.0
(2009). Datos facilitados por Seward Rutkove. DOI:
\href{https://doi.org/10.13026/C24S3D}{10.13026/C24S3D}.
\href{https://physionet.org/content/emgdb/1.0.0/}{Dataset} y
\href{https://physionet.org/content/emgdb/view-license/1.0.0/}{licencia
ODC-By 1.0}.

{[}2{]} PhysioNet.
\href{https://physionet.org/files/emgdb/1.0.0/emg_neuropathy.hea}{Cabecera
del registro de neuropatía}. Fuente de muestreo, calibración y etiqueta
clínica.

{[}3{]} PhysioNet.
\href{https://physionet.org/files/emgdb/1.0.0/emg_healthy.hea}{Cabecera
del registro sin historia de enfermedad neuromuscular}.

{[}4{]} PhysioNet.
\href{https://physionet.org/files/emgdb/1.0.0/emg_myopathy.hea}{Cabecera
del registro de miopatía}.

{[}5{]} PhysioNet.
\href{https://physionet.org/files/emgdb/1.0.0/SHA256SUMS.txt}{SHA256SUMS
de EMGDB 1.0.0}. Base de las comprobaciones de integridad.

{[}6{]} Nandedkar, S. D., Sanders, D. B. y Stålberg, E. V. (1988).
\emph{EMG of reinnervated motor units: a simulation study}.
Electroencephalography and Clinical Neurophysiology, 70(2), 177--184.
DOI:
\href{https://doi.org/10.1016/0013-4694(88)90117-4}{10.1016/0013-4694(88)90117-4}.
Referencia sobre relación entre organización de unidades y medición; no
valida los indicadores de este prototipo.

{[}7{]} SciPy. Documentación oficial de
\href{https://docs.scipy.org/doc/scipy/reference/generated/scipy.signal.welch.html}{Welch}
y
\href{https://docs.scipy.org/doc/scipy/reference/generated/scipy.signal.sosfiltfilt.html}{filtrado
SOS bidireccional}.

{[}8{]} Streamlit.
\href{https://docs.streamlit.io/develop/concepts/architecture}{Modelo de
ejecución} y
\href{https://docs.streamlit.io/develop/api-reference/charts/st.plotly_chart}{gráficos
Plotly}.

{[}9{]} Quarto.
\href{https://quarto.org/docs/output-formats/pdf-basics.html}{PDF
Basics}.

{[}10{]} Escuela Colombiana de Ingeniería Julio Garavito. \emph{Señales
Bioeléctricas y Visualización Interactiva}. Microcurrículo adjunto, sin
vigencia diligenciada. Referencia curricular, no fuente de criterios
diagnósticos.

\textbf{Cambios y atribución del producto derivado.} Se conservaron los
archivos de origen; el software convierte cuentas a microvoltios, filtra
y calcula descriptores. Las bandas, ventanas, tolerancias y diseño de
interfaz son decisiones docentes de esta guía. Los autores del dataset
no avalan esta implementación. No se produjeron resultados
longitudinales ni nuevos diagnósticos.




\end{document}
